\begin{center}
\textsc{Gauss} an \textsc{Bessel}.\quad Göttingen, 1.~April 1827.
\end{center}

\par\noindent\hspace*{\parindent}.\ .\ .\ .\ .\ .\quad Ich denke in diesem Frühjahr die Amplitudo des Bogens zwischen
Göttingen und Altona mit dem Zenithsector zu messen und bin selbst begierig
auf das Resultat. Eine Reihe Beobachtungen, die an den Meridiankreisen Anfangs
1824 gemacht war, gab \(2^{0}\,1^{\prime}\,58^{\prime\prime}\); die Polhöhen, wie ich die meinige nach den
besten Beobachtungen annehme und wie \textsc{Schumacher} die seinige angibt, geben
eher 1 oder \(1\frac{1}{2}\) Secunden weniger; die geodätischen Messungen hingegen
unter Voraussetzung der gleichförmigen Gestalt der Erde und \textsc{Walbecks} Dimen-
\par\clearpage
\noindent{\small\textbf{ZUR HANNOVERSCHEN TRIANGULATION.}}\hfill363\par\medskip
\noindent sionen geben 4 oder 5 Secunden mehr, und fast genau ebenso viel, wie dies
letztere Resultat, geben ganz gleichzeitige Zenithdistanzen von Zenithalsternen,
die ich hier und \textsc{Nehus} in Altona am Meridiankreise beobachtet haben, indem
der Nullpunkt mit Collimatoren bestimmt war. \textsc{Nehus} ist jetzt hier, mir bei
den Sectorbeobachtungen zu helfen; er hat einen \textsc{Repsolds}chen pensilen Colli\-%
mator mitgebracht, bis jetzt aber halte ich die Methode des Nadirpunkts durch
Quecksilberreflexion für bedeutend genauer; doch, denke ich, wird jene Methode
sich noch verbessern lassen.\quad.\ .\ .\ .\ .\ .
